\documentclass[10pt,journal,onecolumn]{IEEEtran}

\usepackage[T1]{fontenc}
\usepackage[utf8]{inputenc}
\usepackage{microtype}
\usepackage{amsmath,amssymb,amsthm,mathtools,bm}
\usepackage{booktabs,tabularx,array,multirow,longtable}
\usepackage{algorithm}
\usepackage[noend]{algpseudocode}
\usepackage{enumitem}
\usepackage{graphicx}
\usepackage{cite}
\usepackage{xurl}
\usepackage[hidelinks]{hyperref}
\usepackage{xcolor}
\usepackage{siunitx}
\usepackage{pdflscape}

\interdisplaylinepenalty=2500
\allowdisplaybreaks
\emergencystretch=2em
\setlength{\LTpre}{3pt}
\setlength{\LTpost}{3pt}

\newtheorem{assumption}{Assumption}[section]
\newtheorem{definition}{Definition}[section]
\newtheorem{lemma}{Lemma}[section]
\newtheorem{proposition}{Proposition}[section]
\newtheorem{theorem}{Theorem}[section]
\newtheorem{corollary}{Corollary}[section]
\newtheorem{remark}{Remark}[section]

\DeclareMathOperator*{\argmin}{arg\,min}
\DeclareMathOperator*{\argmax}{arg\,max}
\newcommand{\R}{\mathbb{R}}
\newcommand{\one}{\mathbf{1}}
\newcommand{\Mone}{\ensuremath{\mathrm{M1}}}
\newcommand{\MoneLP}{\ensuremath{\mathrm{M1\text{-}LP}}}
\newcommand{\Mtwo}{\ensuremath{\mathrm{M2}}}
\newcommand{\Mthree}{\ensuremath{\mathrm{M3}}}
\newcommand{\MtwoMatched}{\ensuremath{\mathrm{M2\text{-}MATCHED}}}
\newcommand{\lexleq}{\preceq_{\mathrm{lex}}}
\newcommand{\pos}[1]{\left[#1\right]_+}
\newcommand{\norm}[1]{\left\lVert#1\right\rVert}
\newcommand{\abs}[1]{\left|#1\right|}
\newcommand{\eps}{\varepsilon}

\title{Technical Appendix and Data Report for\\
Annual Entitlement Accounting for Partial Restrictions in Non-Firm
Transmission Access: A Dutch Transmission Study}

\author{Alban Kryeziu, Stuart X.~Zhu, Tolga Karabas, and Evrim Ursavas%
\thanks{The authors are with the Department of Operations, Faculty of Economics
and Business, University of Groningen, Groningen, The Netherlands.}}

\begin{document}
\maketitle

\begin{center}
\small
Submission-companion edition 2.0 --- 3 September 2026\\
\textit{This edition fixes the scientific, data, computational, and evidence
contract. Numerical outcome tables are inserted only from the final claim-gated
artifact registry.}
\end{center}

\tableofcontents
\listoftables
\listofalgorithms

\section{Document Status, Scope, and Reading Guide}
\label{sec:report_status}

This report is the technical companion to the IEEE Transactions on Power
Systems article. The article states the common operator problem, the three debit
rules, the central set-inclusion certificate, the state-conditioned extension,
and the principal Dutch evidence. This report records the material needed to
audit those statements: data origins and transformations, network and device
schemas, background-state construction, participant mapping, outage semantics,
complete proofs, solver certificates, selected-state engineering tests, stress
protocols, software interfaces, and claim-to-artifact gates.

The study follows one fixed direction,
\begin{equation}
\text{touched support}
\longrightarrow
\text{proportional physical depth}
\longrightarrow
\text{state-conditioned proportional depth}.
\label{eq:report_story}
\end{equation}
The physical action is always implemented restriction. M3 learns a debit-rate
card, not the action. The parent annual and rolling problems contain restriction
only. Common redispatch, independent DC reconstruction, and AC power flow are
separate evidence layers applied after policy decisions are frozen.

The report distinguishes three factual statuses. Method statements are fixed by
this document. Data statements are admitted only when the corresponding source,
unit, time-grid, mapping, and hash records exist. Numerical outcome statements
are admitted only when the final campaign and claim gate identify their exact
supporting artifacts. A missing data table or failed gate is not completed by
prose. It remains a listed submission blocker until the machine-readable file is
supplied.

\section{Notation and Frozen Scientific Contract}
\label{sec:report_contract}

\subsection{Sets, indices, and time}

Let $y\in\mathcal Y$ index years, $t\in\mathcal T_y$ intervals,
$n\in\mathcal N$ directional non-firm positions, $b\in\mathcal B$ buses,
$c\in\mathcal C$ network states, and $r\in\mathcal R_c$ monitored elements in
state $c$. The native final grid is hourly, but each interval carries a
positive duration $\Delta_t$ in hours. Timestamps are unique, strictly
ordered, timezone-aware, and complete for the declared year. No quarter-hour
variation is synthesized from hourly data.

Each position has contracted capacity $P_n>0$, request
$y^{\mathrm{req}}_{n,t}\geq0$, contract-duration indicator
$\kappa_{n,t}\in\{0,1\}$, and restriction-eligibility indicator
$\chi_{n,t}\in\{0,1\}$. The consistency rule is
$\chi_{n,t}\leq\kappa_{n,t}$. Positive request, contractual duration, and
restriction eligibility are never inferred from one another.

\subsection{Frozen methodological decisions}

The following decisions are normative.
\begin{enumerate}[leftmargin=*,itemsep=2pt]
\item The parent annual and rolling model uses restriction only. Common
redispatch is a separate selected-state audit.
\item Same-scale M1--M3 cases use the same action bounds, requests, network
evaluator, ratings, participant set, objective, and actual allowance vector.
\item Annual allowance is computed from the explicit contract-duration mask
using floor rounding on the native hourly grid.
\item M1 retains binary activation and a two-sided materiality link.
\item The reported operator value is the exact Stage-I/Stage-II pair $(S,E)$.
Stage III is optional reporting-only selection on that face. It is omitted from
the primary annual and rolling campaigns and never constructs an M3 teacher.
\item M3 teacher multipliers and primal actions come from the same exact,
unperturbed Stage-II M2 solve.
\item Every deployable time-varying feature passes the 07:00
Europe/Amsterdam D--1 availability gate. Scenario or realized residual load is
not relabeled as a forecast.
\item The final deployable card uses F0 only: deterministic calendar variables,
contracted capacity, request fraction, zone, and technology. Direction defines
the ranking class.
\item A direction with fewer than two positive-reference-FDE buckets receives
$q=1$.
\item M2-MATCHED is an ex-post participant-specific allowance experiment, not a
scalar-$\rho$ policy.
\item Background states retain separate demand, PV, onshore-wind, and offshore-
wind components and use one frozen balancing convention.
\item The detailed 383-bus layer supports localization and validation; the
45-bus EHV layer supports policy optimization.
\item The final primary population is the 79 directly represented positions.
The 35 multibus positions remain sensitivity-only because the complete all-or-
none embedding gate is not passed.
\item The final evaluator contains 95 monitored elements, the intact state, and
93 retained supported non-islanding outages, for 8,837 state--element rows per
hour.
\item The outage claim is limited to identifiable retained single modeled
elements. Parallel-equivalent and phase-shift semantics are audited explicitly.
\item The final implementation is a clean V3 core consuming immutable experiment
objects. Legacy scripts are computational and provenance references only.
\end{enumerate}

\begin{table}[t]
\centering
\caption{Frozen primary physical and chronological case.}
\label{tab:report_primary_case}
\begin{tabularx}{\textwidth}{@{}lX@{}}
\toprule
Object & Final specification \\
\midrule
Time grid & 61,344 hourly intervals, 2025--2031; timezone-aware and complete \\
Policy network & 45 retained 220/380-kV buses; 95 monitored line/transformer elements \\
Network states & intact plus 93 retained non-islanding outages \\
Row universe & 8,837 state--element rows per hour, with positive and negative limits \\
Participants & 79 direct primary; 35 multibus sensitivity-only \\
Card split & 2025--2027 fit; 2028 validation; 2029--2030 held out; 2031 confirmation \\
Primary allowances & $\rho=0.15$ for M1--M2; $\rho=0.03$ for the M3 mechanism study \\
Rolling transfer & 48-h look-ahead; 24-h implementation; no terminal reserve or terminal value \\
\bottomrule
\end{tabularx}
\end{table}

\subsection{Immutable experiment identity}

Every solve is represented by an immutable experiment object. At minimum it
contains: case ID, physical-case ID and hash, time-grid ID and hash, rule,
nominal allowance share, allowance metric, allowance source, actual allowance
vector and hash, contract-mask and eligibility-mask hashes, participant-set
ID, contingency-set ID and hash, card and feature-set IDs, information level,
look-ahead and implementation durations, solver profile, source-code hash,
and units. The solver refuses an object with a missing or inconsistent
scientific field.

\section{Common Security-Constrained Operating Model}
\label{sec:report_common_model}

\subsection{Admissible action and materiality}

The restriction bound is
\begin{equation}
0\leq x_{n,t}\leq\bar x_{n,t}
:=\chi_{n,t}\min\{P_n,y^{\mathrm{req}}_{n,t}\}.
\label{eq:report_xbar}
\end{equation}
The active-coordinate set is
$\mathcal I_y=\{(n,t):\bar x_{n,t}>0\}$. For M1, the minimum implementable
restriction is
\begin{equation}
\underline x_{n,t}
:=\min\left\{\bar x_{n,t},\vartheta_n\right\},
\qquad
\vartheta_n:=\max\{1\ \mathrm{MW},0.01P_n\}.
\label{eq:report_floor}
\end{equation}
The threshold defines the accounting action, not a measurement tolerance. A
separate numerical tolerance is used only to replay zero and integrality.

\subsection{Balanced background and prospective transactions}

Let $\bm p_t^{\mathrm{bg}}\in\R^{|\mathcal B|}$ be the frozen background active
injection before prospective non-firm transactions. Let
$\bm A=[\bm a_n]_{n\in\mathcal N}$ be the participant-transaction matrix, with
$-\bm a_n$ the nodal effect of one MW of requested transport and
$+\bm a_n$ the effect of restricting one MW. The mandatory balance identities
are
\begin{equation}
\one^\top\bm p_t^{\mathrm{bg}}=0,
\qquad
\one^\top\bm a_n=0\quad\forall n.
\label{eq:report_balance}
\end{equation}
The requested and post-restriction states are
\begin{align}
\bm p_t^{\mathrm{req}}
&=\bm p_t^{\mathrm{bg}}-\bm A\bm y_t^{\mathrm{req}},
\notag\\
\bm p_t(\bm x)
&=\bm p_t^{\mathrm{req}}+\bm A\bm x_t.
\label{eq:report_states}
\end{align}
All columns of $\bm A$ use the same frozen counter-injection convention and
carry the balancing-key hash used to build them.

\subsection{Affine DC evaluator and thermal rows}

For each network state $c$, the DC equations are represented in affine form
\begin{equation}
\bm f_t^c=\bm H^c\bm p_t(\bm x)+\bm h^c.
\label{eq:report_affine}
\end{equation}
The offset $\bm h^c$ accounts for phase-shifting transformers and any fixed
branch-flow term arising when voltage-angle equations are eliminated. It is
serialized as a zero vector only after the device audit proves that all
relevant fixed terms vanish.

For row $r\in\mathcal R_c$ with rating $\bar F_r^c>0$,
\begin{equation}
-\bar F_r^c-s_{r,t,c}^{-}
\leq f_{r,t}^c
\leq\bar F_r^c+s_{r,t,c}^{+},
\qquad s_{r,t,c}^{\pm}\geq0.
\label{eq:report_rows}
\end{equation}
The row ID is a composite of network-state ID, physical-element ID,
orientation, and rating source. Array shape is never accepted as a substitute
for exact ID alignment.

\subsection{Sequential operator objective}

Stage I minimizes
\begin{equation}
S(\bm x,\bm s)
=\sum_{t\in\mathcal T_y}\Delta_t
\sum_{c\in\mathcal C}\sum_{r\in\mathcal R_c}
(s_{r,t,c}^{+}+s_{r,t,c}^{-}).
\label{eq:report_S}
\end{equation}
Stage II fixes
$S\leq S^{\star}+\eps_S$ and minimizes
\begin{equation}
E(\bm x)=\sum_{n,t}\Delta_tx_{n,t}.
\label{eq:report_E}
\end{equation}
The reported value is $\mathcal V=(S^{\star},E^{\star})$.

A downstream figure or engineering manifest may require one selected policy
vector. In that case, an optional Stage III fixes
$S\leq S^{\star}+\eps_S$ and $E\leq E^{\star}+\eps_E$ and applies a declared
deterministic selector. This selector has no scientific objective status and is
not required for the primary annual or rolling evidence. The final production
campaign omits it unless the selected vector itself is an explicit artifact.
Any Stage-III result records its objective, tolerances, and source hashes and
must reproduce the Stage-I/II values. Stage III is never used to construct the
M3 teacher.

\subsection{Units and interpretation}

$S$ is measured in MW$\cdot$h$\cdot$row-state and is an integrated residual
thermal-row index. $E$ is measured in MWh and equals total restricted energy.
Neither quantity is a price, welfare measure, customer cost, actual saving, or
reinforcement value. A zero or lower DC value does not establish AC security.
Baseline background and requested-state violations are stored separately so
that public-data background insecurity is not silently attributed to a debit
rule.

\section{Annual Entitlement Geometry}
\label{sec:report_geometry}

\subsection{Allowance construction}

On the uniform grid,
\begin{equation}
N_{n,y}^{\mathrm{ctr}}=\sum_t\kappa_{n,t},
\qquad
K_{n,y}(\rho)=\lfloor\rho N_{n,y}^{\mathrm{ctr}}\rfloor,
\qquad
H_{n,y}(\rho)=K_{n,y}(\rho)\Delta.
\label{eq:report_allowance}
\end{equation}
The resulting duration vector is serialized. Same-scale cases consume the
same bytes, not merely the same nominal $\rho$. On a nonuniform grid, $H$ is
an exogenous duration allowance and no integer-interval interpretation is
imposed.

\subsection{M1, M1-LP, M2, and M3}

M1 uses
\begin{align}
\underline x_{n,t}u_{n,t}
&\leq x_{n,t}\leq\bar x_{n,t}u_{n,t},
\qquad u_{n,t}\in\{0,1\},
\label{eq:report_m1_links}\\
D_n^1(\bm u)
&=\sum_t\Delta_tu_{n,t}\leq H_n.
\label{eq:report_m1}
\end{align}
M2 uses
\begin{equation}
D_n^2(\bm x)=\sum_t\Delta_t\frac{x_{n,t}}{P_n}\leq H_n.
\label{eq:report_m2}
\end{equation}
M3 uses a frozen positive card,
\begin{equation}
D_n^3(\bm x;\bm q)
=\sum_t\Delta_tq_{n,t}\frac{x_{n,t}}{P_n}\leq H_n.
\label{eq:report_m3}
\end{equation}
Equal numerical allowances do not imply equal contractual services. M1 limits
materially touched time. M2 limits cumulative normalized physical exposure.
M3 changes the state-dependent debit exchange rate applied to that exposure.

\subsection{Characterization of proportional depth}

\begin{proposition}[Normalized additive physical-exposure debit]
\label{prop:unique_m2}
For fixed $P>0$, let $d(\cdot;P):[0,P]\to\R_+$ satisfy
$d(0;P)=0$, $d(P;P)=1$, monotonicity, and
$d(x_1+x_2;P)=d(x_1;P)+d(x_2;P)$ whenever
$x_1,x_2,x_1+x_2\in[0,P]$. Then $d(x;P)=x/P$.
\end{proposition}

\begin{IEEEproof}
For rational $r=k/m\in[0,1]$, additivity gives
$m d(P/m;P)=d(P;P)=1$ and hence $d(rP;P)=r$. For arbitrary
$x\in[0,P]$, choose rational sequences $r_j\uparrow x/P$ and
$s_j\downarrow x/P$. Monotonicity gives
$r_j\leq d(x;P)\leq s_j$; taking limits yields $d(x;P)=x/P$.
\end{IEEEproof}

The result characterizes a normalized physical-exposure unit. It does not
assert that touched duration and physical exposure encode the same burden or
should receive the same legal treatment.

\subsection{Exact M1-LP projection and isolated convex hull}

For one participant, define the isolated M1 set
\begin{equation}
\mathcal U_n
=\left\{(\bm x,\bm u):
\eqref{eq:report_m1_links},\ \eqref{eq:report_m1}\right\}.
\end{equation}

\begin{theorem}[Projection, hierarchy, and strict certificate]
\label{thm:report_hierarchy}
The projection of the continuous relaxation of $\mathcal U_n$ is
\begin{equation}
0\leq x_{n,t}\leq\bar x_{n,t},
\qquad
D_n^{1L}(\bm x)
:=\sum_{t:\bar x_{n,t}>0}\Delta_t
\frac{x_{n,t}}{\bar x_{n,t}}
\leq H_n.
\label{eq:report_m1lp}
\end{equation}
If $\Delta_t=\Delta$ and $H_n=K_n\Delta$ with integer $K_n$, this projected
polytope equals $\operatorname{conv}(\operatorname{proj}_x\mathcal U_n)$.
The equality is only participant-level and precedes intersection with the
network constraints.

Moreover,
\begin{equation}
D_n^{1L}(\bm x)-D_n^2(\bm x)
=\sum_{t:\bar x_{n,t}>0}\Delta_tx_{n,t}
\left(\frac1{\bar x_{n,t}}-\frac1{P_n}\right)\geq0,
\label{eq:report_gap}
\end{equation}
and, under any common physical model and allowance vector,
\begin{equation}
\Pi\mathcal F^{\Mone}(\bm H)
\subseteq\mathcal F^{\MoneLP}(\bm H)
\subseteq\mathcal F^{\Mtwo}(\bm H).
\label{eq:report_hierarchy}
\end{equation}
Consequently,
$\mathcal V^{\Mtwo}\lexleq\mathcal V^{\MoneLP}
\lexleq\mathcal V^{\Mone}$.

Let $Q$ be M2 or a fixed-card M3 case with the same physical model and
allowance. If $\overline S^Q$ is replay-feasible,
$\underline S^{1L}$ is a certified M1-LP lower bound, and
\begin{equation}
\overline S^Q+\eps_{\mathrm{cert}}<\underline S^{1L},
\label{eq:report_strict}
\end{equation}
then $S^{Q,\star}<S^{\Mone,\star}$.
\end{theorem}

\begin{IEEEproof}
For $\bar x_{n,t}>0$, the relaxed upper link requires
$u_{n,t}\geq x_{n,t}/\bar x_{n,t}$. Setting
$u_{n,t}=x_{n,t}/\bar x_{n,t}$ satisfies the lower link because
$\underline x_{n,t}\leq\bar x_{n,t}$ and produces
\eqref{eq:report_m1lp}; no other projected restriction is needed.

On the uniform grid, put $z_t=x_{n,t}/\bar x_{n,t}$. The polytope
$0\leq z_t\leq1$, $\sum_tz_t\leq K_n$ is the independence polytope of a
uniform matroid. Its extreme points are binary. Mapping a binary extreme point
back by $x_{n,t}=\bar x_{n,t}z_t$ gives an M1-feasible projected point because
$\underline x_{n,t}\leq\bar x_{n,t}$. This proves the isolated convex-hull
claim. In general,
$\operatorname{conv}(X\cap Y)\neq\operatorname{conv}(X)\cap Y$, so no
network-intersected equality is asserted.

Equation~\eqref{eq:report_gap} follows by subtraction and is nonnegative since
$\bar x_{n,t}\leq P_n$. Relaxing integrality gives the first inclusion in
\eqref{eq:report_hierarchy}; the debit inequality gives the second. Nested
feasible sets give the lexicographic value ordering. Finally,
$S^{Q,\star}\leq\overline S^Q$ and
$\underline S^{1L}\leq S^{\MoneLP,\star}\leq S^{\Mone,\star}$, so
\eqref{eq:report_strict} implies the strict result.
\end{IEEEproof}

\subsection{Rolling upper endpoint}

Let the implemented blocks $\mathcal I_\tau$ form a disjoint partition of the
annual chronology. For M1,
\begin{equation}
B_{n,\tau+1}=B_{n,\tau}
-\sum_{t\in\mathcal I_\tau}\Delta_tu_{n,t}.
\label{eq:report_m1_ledger}
\end{equation}
If every block is physically feasible and the terminal balances are
nonnegative, telescoping gives an annual M1-feasible trajectory. Hence
\begin{equation}
S^{\Mtwo,\star}\leq S^{\MoneLP,\star}
\leq S^{\Mone,\star}\leq S^{\Mone,\mathrm{roll}}.
\label{eq:report_scalar_bracket}
\end{equation}
This is an annual scalar bracket. Hourly cumulative M1-LP and rolling-M1 curves
show when the fixed trajectories accumulate residual violation; they are not
pointwise optimality bounds.

\subsection{Temporal aggregation}

Let a coarse interval of duration $\Delta$ contain subintervals $j$ with
durations $\delta_j$ and actions $x_j$. With
$\bar x=\Delta^{-1}\sum_j\delta_jx_j$,
\begin{equation}
\Delta\frac{\bar x}{P}
=\sum_j\delta_j\frac{x_j}{P},
\end{equation}
so M2 is average-preserving. M1 is not generally average-preserving when a
heterogeneous fine-grid action is replaced by its average. Four quarter-hours
$(P,0,0,0)$ have fine-grid M1 debit $0.25$ h, whereas their hourly average
$P/4$ has coarse-grid M1 debit $1$ h; both have M2 debit $0.25$ h.

M3 preserves debit under aggregation only when the coarse rate is
exposure-weighted:
\begin{equation}
\bar q
=\frac{\sum_j\delta_jq_jx_j}{\sum_j\delta_jx_j}
\quad\Longrightarrow\quad
\Delta\bar q\frac{\bar x}{P}
=\sum_j\delta_jq_j\frac{x_j}{P}.
\end{equation}
Independent averaging of $q$ and $x$ generally fails this identity.

\section{Exact Screening, Closure, Replay, and Numerical Certification}
\label{sec:report_certification}

\subsection{Action-box row screening}

Let $\bm G^c=\bm H^c\bm A$ and
\begin{equation}
\bm f_t^{c,\mathrm{req}}
=\bm H^c(\bm p_t^{\mathrm{bg}}-\bm A\bm y_t^{\mathrm{req}})+\bm h^c.
\end{equation}
For row $(c,r,t)$, the restriction box
$0\leq x_{n,t}\leq\bar x_{n,t}$ implies
\begin{align}
L_{r,t}^c
&=f_{r,t}^{c,\mathrm{req}}
+\sum_n\min\{0,G_{r,n}^c\bar x_{n,t}\},
\notag\\
U_{r,t}^c
&=f_{r,t}^{c,\mathrm{req}}
+\sum_n\max\{0,G_{r,n}^c\bar x_{n,t}\}.
\label{eq:report_screen_bounds}
\end{align}

\begin{lemma}[Safety of action-box screening]
If $L_{r,t}^c\geq-\bar F_r^c$ and
$U_{r,t}^c\leq\bar F_r^c$, row $(c,r,t)$ cannot be violated by any admissible
restriction vector and may be omitted from the initial model.
\end{lemma}

\begin{IEEEproof}
The linear row effect is separable over the box. Its minimum is obtained by
setting $x_{n,t}=\bar x_{n,t}$ when $G_{r,n}^c<0$ and zero otherwise; its
maximum uses the opposite choice. These extrema are exactly
\eqref{eq:report_screen_bounds}. If both lie within the rating interval, every
point in the box lies within it.
\end{IEEEproof}

The screening rule is independent of the debit law and is rerun only when the
physical case or action box changes. The final scientific implementation does
not remove a participant-time action merely because it is absent from one
legacy screened model; any action pruning must have an equally explicit
dominance proof and a dedicated regression test.

\subsection{Potential rows, initial seed, and exact constraint generation}

The action-box test has two uses that must not be conflated. Its full potentially
active set is used to prove whether a participant--hour action can be removed by
dominance. The first optimization model may use a much smaller deterministic
row seed containing requested-state violations and near-binding rows. Every
registered row remains in the closure universe.

\begin{theorem}[Exact seeded row generation]
\label{thm:seeded_closure}
Consider one lexicographic stage. Let the reduced model contain any subset of the
registered directed thermal rows and the corresponding slack terms. Solve the
reduced stage to its required certificate, evaluate every omitted row at the
candidate, and add every omitted row with positive required slack above the
closure tolerance. If the procedure terminates with zero omitted-row slack and
the candidate passes full replay, then the candidate objective equals the
objective of the complete-row stage.
\end{theorem}

\begin{IEEEproof}
Removing rows and their nonnegative slack terms gives a lower bound on the
complete objective. At termination, the reduced optimum is feasible for the
complete model by assigning zero slack to every omitted row. Its complete and
reduced objectives therefore coincide. The complete optimum is at most this
value and at least the reduced optimum, so all three values are equal. The Stage-I optimizer itself is evaluated on the complete row universe before
its objective is accepted, so an omitted row cannot make the Stage-I reference
artificially low. The same row-closure argument applies to Stage II after the
certified Stage-I cap is included. If Stage-II closure adds a row that can alter
the Stage-I optimum, Stage I is reclosed before Stage II is accepted.
\end{IEEEproof}

Near-binding rows may be added with a deterministic margin to reduce closure
rounds. Such additions restrict neither the declared full model nor the proof;
they merely enlarge the seed. A cached row set from another rule or year may be
used only as an additional seed. It never licenses omission from final replay.

\subsection{Iterative full-row closure}

Screening is followed by exact closure rather than trusted as a terminal
approximation. Algorithm~\ref{alg:closure} applies separately to Stage I and
Stage II.

\begin{algorithm}[t]
\caption{Signed-row closure for one lexicographic stage}
\label{alg:closure}
\begin{algorithmic}[1]
\Require immutable physical case, initial safely retained rows, rule, allowance
\Repeat
  \State solve the current stage to its required certificate
  \State reconstruct $f_{r,t}^c$ for every registered $(c,r,t)$ outside the model
  \State identify omitted positive and negative violations above $\eps_{\rm close}$
  \If{no hidden violation exists}
    \State \Return solution and complete row manifest
  \EndIf
  \State add every hidden directed row and preserve all previous rows
\Until{closure}
\end{algorithmic}
\end{algorithm}

If Stage-II closure adds rows, Stage I is reclosed because the Stage-I optimum
may change. In addition, Stage II is checked against the independently
reconstructed aggregate physical face. Write
\begin{equation}
S(\bm x)=\sum_{t,c,r}\Delta_t\left[\left|f_{r,t}^{c,\rm req}+\bm G_r^c\bm x_t\right|-\bar F_r^c\right]_+.
\label{eq:report_physical_S_polyhedral}
\end{equation}
This function is convex and polyhedral. If a Stage-II candidate $\bm x^k$ has
$S(\bm x^k)>S^{\rm cap}+\eps_{\rm face}$, the identical mathematical model is first re-solved with the numerical recovery profile. If the replay violation persists, choose any $\bm g^k\in\partial S(\bm x^k)$ and append
\begin{equation}
S(\bm x^k)+(\bm g^k)^\top(\bm x-\bm x^k)\leq S^{\rm cap}.
\label{eq:report_physical_face_cut}
\end{equation}
For every physically admissible point with $S(\bm x)\leq S^{\rm cap}$,
convexity implies the left-hand side of \eqref{eq:report_physical_face_cut} is
at most $S(\bm x)$ and hence at most the cap. The cut is therefore globally
valid and removes the offending point without removing any point on the true
Stage-I face. The Stage-II objective is then re-solved. Since $S$ has finitely
many affine pieces, repeated separation is finite in exact arithmetic; the
implementation imposes a finite round limit and fails closed if it is reached.
No post-solve projection or residual centering is permitted. The optional
Stage-III selector is subjected to the same physical-face replay and separation
logic.

The closure log records every generation, added row ID, physical-face cut,
violation, objective, solver status, and source hash. A stalled closure, hidden
row, or unresolved aggregate face violation after the permitted rounds is a
hard failure, not a warning.

\subsection{Independent replay}

Replay is implemented in a module that receives only immutable artifacts and
solution arrays. It does not import the solver model or reuse solver
constraint expressions. It reconstructs:
\begin{enumerate}[leftmargin=*,itemsep=2pt]
\item $\bar x$ and the action-bound residuals;
\item M1 activation and materiality consistency when applicable;
\item every participant debit and allowance residual;
\item every rolling balance transition;
\item $\bm p_t^{\mathrm{req}}$, $\bm p_t(\bm x)$, and every affine row flow;
\item $S$, $E$, touched duration, physical FDE, and weighted debit;
\item exact experiment, physical-case, row, participant, and card hashes.
\end{enumerate}
A solution is replay-feasible only when all residuals are within their
predeclared units and tolerances. The replayed value, not the solver-reported
objective alone, is used in the strict M1 certificate.

\subsection{Tolerance and certificate contract}

Table~\ref{tab:core_tolerances} gives the default numerical contract. A change
is a methodology change and requires a new freeze version.

\begin{table}[t]
\centering
\caption{Core numerical tolerances.}
\label{tab:core_tolerances}
\begin{tabular}{@{}lll@{}}
\toprule
Quantity & Default & Role \\
\midrule
Primary solver feasibility & $10^{-8}$ & primal feasibility target \\
Primary solver optimality & $10^{-8}$ & dual/optimality target \\
Generic physical-face comparison & $\max\{10^{-7},10\,\texttt{FeasibilityTol}\}$ & trigger replay separation \\
Annual support-OA face comparison & $\max\{10^{-6},100\,\texttt{FeasibilityTol}\}$ & support-oracle contract \\
Recovery feasibility/optimality & $10^{-9}$ & replay-triggered generic re-solve only \\
Lexicographic absolute tolerance & $10^{-6}$ & scientific Stage-I/II face \\
Lexicographic relative tolerance & $10^{-9}$ & scale-adjusted scientific face \\
Touch-zero tolerance & $10^{-7}$ MW & numerical zero only \\
Selected-dual tolerance & $2\times10^{-6}$ & sign and repeat audit \\
AC ordering tie & $10^{-4}$ & normalized thermal-exceedance comparison \\
\bottomrule
\end{tabular}
\end{table}

The two lexicographic tolerances in the table define the scientific optimization problem. They are not enlarged when a solver produces a numerically inconsistent point. The generic numerical face comparison is only an acceptance/recovery trigger. When it fails, the physical support cut \eqref{eq:report_physical_face_cut} is added and the same scientific problem is re-solved. A recovery solve may set \texttt{NumericFocus=3}, \texttt{Quad=1}, and $10^{-9}$ feasibility/optimality tolerances; these settings alter solver arithmetic only. The returned vector is never edited after optimization.

For a candidate $Q$, define the replay-safe feasible value
\begin{equation}
\overline S^Q
=S^{Q,\mathrm{replay}}
+\eps_{\mathrm{sum}}^Q,
\end{equation}
where $\eps_{\mathrm{sum}}^Q$ is the reported floating-point accumulation
margin. Let the LP lower bound be
\begin{equation}
\underline S^{1L}=B^{1L}_{\mathrm{solver}}-\eps_{\mathrm{LB}}.
\end{equation}
The replay and lower-bound corrections are already included in
$\overline S^Q$ and $\underline S^{1L}$. The additional strict-comparison
margin is therefore
\begin{equation}
\eps_{\mathrm{cert}}
=\max\{10^{-6},10^{-9}\max(1,\abs{\overline S^Q},\abs{\underline S^{1L}})\}.
\label{eq:report_eps_cert}
\end{equation}
The certificate passes only through the strict inequality
\eqref{eq:report_strict}. No percentage is reported when that inequality does
not hold.

\subsection{Continuous and mixed-integer statuses}

M1-LP, M2, M3, and M2-MATCHED are continuous problems. Their primary statuses
must be optimal and their dual/bound artifacts must be available. Rolling M1
is a MILP and must supply a feasible implemented trajectory with certified
per-epoch status under the declared absolute-gap policy. A time-limited
annual-M1 incumbent is not used as a scientific comparator. The annual M1
lower and upper endpoints come from M1-LP and the concatenated rolling-M1
trajectory, respectively.

\section{M3 Teacher, Prediction, and Frozen Rate Card}
\label{sec:report_m3}

\subsection{Canonical teacher problem}

The teacher table, target, candidate cards, validation results, and frozen card
are allowance-specification specific. The nominal share, actual allowance vector,
contract-mask hash, and allowance-vector hash are part of the card ID. A card
built under one allowance specification is never reused at another share as the
same scientific object. The allowance index is suppressed below only to reduce
notation.

For $y\in\{2025,2026,2027,2028\}$, the teacher is the M2 problem with the
exact Stage-I and Stage-II objectives in
\eqref{eq:report_S}--\eqref{eq:report_E}. Let $\mu_{n,y}\geq0$ be the
multiplier of $D_n^2(\bm x)\leq H_{n,y}$ in the selected Stage-II primal--dual
pair. The attribution is
\begin{equation}
\alpha_{n,t,y}
=\mu_{n,y}\Delta_t\frac{x_{n,t,y}^{\Mtwo}}{P_n}.
\label{eq:report_alpha}
\end{equation}

For a fixed teacher Stage-I cap
$\bar S_y=S_y^{\Mtwo,\star}+\varepsilon_S$, define the conditional
parametric value
\begin{equation}
\varphi_y(\bm q;\bar S_y)
=\min\left\{E(\bm x):\ (\bm x,\bm s)\in\mathcal X_y,\ 
S(\bm x,\bm s)\leq\bar S_y,\ 
D_n^3(\bm x;\bm q)\leq H_{n,y}\quad\forall n\right\}.
\label{eq:report_conditional_value}
\end{equation}
At the flat card $\bm q=\one$, if the selected basis is locally stable or
$\varphi_y$ is differentiable, standard parametric-LP sensitivity gives
\cite{BonnansShapiro2000}
\begin{equation}
\frac{\partial\varphi_y}{\partial q_{n,t}}
=\mu_{n,y}\Delta_t\frac{x_{n,t,y}^{\Mtwo}}{P_n}=\alpha_{n,t,y}.
\end{equation}
The Stage-I cap is held fixed in this statement. The attribution is not a
derivative of the recomputed Stage-I optimum or of the complete lexicographic
value. Without the stated regularity, only the algebraic selected
primal--dual attribution is claimed. The report never calls $\alpha$ an
electricity price, congestion price, Stage-I gradient, unique shadow price, or
causal effect.

\subsection{Average-tie percentile target}

For each direction $g\in\{\mathrm{import},\mathrm{export}\}$, let
$\mathcal J_g$ be the fitting rows and $m_g=|\mathcal J_g|$. The target is
\begin{equation}
r_i
=\frac{\#\{j\in\mathcal J_g:\alpha_j<\alpha_i\}
+\tfrac12(1+\#\{j\in\mathcal J_g:\alpha_j=\alpha_i\})}{m_g}.
\label{eq:report_rank_target}
\end{equation}
This is exactly the average rank divided by $m_g$. Ties at zero are retained;
they are not jittered. When the training table exceeds 200,000 rows in one
direction, deterministic direction-stratified sampling with master seed 118
is used for model fitting. Bucket inputs and rate normalization are computed
on the full fitting table rather than only the sample.

\subsection{Information-provenance gate}

The cutoff is 07:00 Europe/Amsterdam on D--1. Every feature row records source
ID, file hash, availability class, timezone, revision policy, and cutoff status.
A forecast is eligible only when an archived vintage existed before the cutoff
and was not revised using later information. A scenario trajectory, realized
quantity, post-event correction, security outcome, KKT multiplier, future
restriction, future error, or held-out outcome is ineligible.

The final deployable family is
\begin{equation}
F_0=\{\text{hour/day calendar variables, capacity, request fraction,
zone, technology}\}.
\label{eq:report_F0}
\end{equation}
Import and export are fitted as separate ranking classes. Static and calendar
variables pass by construction under an explicit availability class; they do
not carry invented forecast issue timestamps. Authenticated provincial and
national residual-load forecasts would define richer families, but those
vintages are not available in the frozen study. Such variables and the former
residual-load public score are excluded from deployable model selection.
Realized residual load is retained only for post-freeze interpretation.

\subsection{Predictor specifications}

The candidate predictor families are intentionally limited.
\begin{enumerate}[leftmargin=*,itemsep=3pt]
\item \emph{Ridge regression.} Numeric variables receive fitting-sample median
imputation and standardization; categorical variables receive one-hot encoding
with unknown-category handling. The ridge parameter is
$\alpha_{\mathrm{ridge}}=2$, distinct from the scarcity label $\alpha$ and the
card regularization $\lambda_q$.
\item \emph{Shallow histogram boosting.} Numeric variables receive fitting-
sample median imputation and standardization; categorical variables receive
ordinal encoding with an explicit unknown value. The frozen profile uses 80
iterations, at most 15 leaves, depth at most 3, minimum leaf size 40,
$\ell_2$ regularization 1, learning rate 0.08, and seed 118.
\end{enumerate}
No additional architecture is introduced after validation begins. The former
public residual-load candidate is not evaluated because its required archived
D--1 vintages fail the information gate.

\subsection{Thresholds, bucket inputs, and rate optimization}

Training-score tertiles are computed separately within direction. Constant
scores produce one bucket and therefore trigger the flat-direction rule.
Otherwise, buckets $b=1,2,3$ are ordered from low to high predicted scarcity.
Define
\begin{align}
A_{g,b}
&=\sum_{i\in(g,b)}\alpha_i,
\notag\\
W_{g,b}^{\mathrm{ref}}
&=\sum_{i\in(g,b)}\Delta_i\frac{x_i^{\Mtwo}}{P_i}.
\label{eq:report_AW}
\end{align}
The primary bounds are $[0.5,1.5]$ and the primary regularization is
$\lambda_q=10^{-3}\,\mathrm{MW}$. Under the definitions in
\eqref{eq:report_AW}, $A_{g,b}$ is measured in MWh and
$W_{g,b}^{\mathrm{ref}}$ in hours, so both terms of the card objective have
units of MWh. The card solves
\begin{align}
\min_{\bm q_g}\quad
&\sum_{b=1}^3A_{g,b}q_{g,b}
+\frac{\lambda_q}{2}\sum_{b=1}^3
 W_{g,b}^{\mathrm{ref}}(q_{g,b}-1)^2
\label{eq:report_card_qp}\\
\mathrm{s.t.}\quad
&q_{g,1}\geq q_{g,2}\geq q_{g,3},
\notag\\
&0.5\leq q_{g,b}\leq1.5,
\notag\\
&\sum_bW_{g,b}^{\mathrm{ref}}q_{g,b}
=\sum_bW_{g,b}^{\mathrm{ref}}.
\notag
\end{align}
A direction with fewer than two positive $W_{g,b}^{\mathrm{ref}}$ values is
set to $(1,1,1)$. If the primary QP has multiple minimizers, its optimal value
is fixed and a secondary QP minimizes $\sum_b(q_{g,b}-1)^2$. The serialized
rate vector must satisfy bounds, order, reference normalization, and KKT
residual tolerances.

The fitting-reference equality is not an exposure-neutrality theorem. Once
M3 changes operational actions, the realized distribution across buckets and
the participant physical FDE may change.

\subsection{Validation selection and freeze}

Predictor parameters, preprocessing, and thresholds are fit using 2025--2027.
Each eligible feature-family/predictor card is then run through the complete
2028 annual M3 operator problem. Selection follows:
\begin{enumerate}[leftmargin=*,itemsep=2pt]
\item minimum validation Stage-I value within the frozen tolerance;
\item minimum validation Stage-II restricted energy within tolerance;
\item minimum standard deviation of the six directional rates;
\item deterministic complexity order: ridge, histogram boosting;
\item stable candidate ID.
\end{enumerate}
A candidate must also pass replay, card KKT, information-provenance, and exact
rebuild gates. If no eligible nonflat candidate passes, the allowance regime
freezes the flat card $q=1$. The selected card and all candidate outcomes are
serialized. A pre-freeze access audit proves that the card-building process did
not read any 2029--2031 feature values, physical inputs, or outcomes through the
fit, threshold, selection, rate, bound, shuffle, or amplitude code path.

The freeze contains: allowance specification and vector hash, feature schema and
provenance hash, training rows and hash, target definition, model family and
parameters, preprocessing, thresholds, rates, rate bounds, $\lambda_q$,
direction fallbacks, software and code hashes, 2028 selection
values, and the held-year access audit. A changed object requires a new card ID
and a new evaluation freeze.

\subsection{Selected-dual and card stability}

The exact-rebuild audit solves each fitting M2 teacher twice with identical
inputs and profile. It requires Stage-I and Stage-II objective agreement under
the lexicographic tolerance, maximum restriction difference at most
$2\times10^{-6}$ MW, maximum selected-dual difference at most
$2\times10^{-6}$, and an identical rebuilt card within the card tolerance.
Failure triggers the flat-card fallback.

Additional profiles vary presolve, LP algorithm, and small Stage-III-only
selection perturbations. These are diagnostics rather than alternate teachers.
They report within-direction Spearman rank agreement, bucket agreement, rate
differences, and whether the 2028 selected candidate changes. The manuscript
uses derivative language only when the corresponding local regularity and
stability evidence is explicitly established.

\subsection{M3 exposure theorem and M2-MATCHED}

\begin{theorem}[Rate bounds and matched exposure]
Suppose $0<q_{\min}\leq q_{n,t}\leq q_{\max}$. Then
\begin{equation}
\mathcal F^{\Mtwo}(\bm H/q_{\max})
\subseteq\mathcal F^{\Mthree}(\bm H,\bm q)
\subseteq\mathcal F^{\Mtwo}(\bm H/q_{\min}).
\label{eq:report_m3_sandwich}
\end{equation}
For $D_n^2(\bm x)>0$,
\begin{equation}
\bar q_n^{\mathrm{use}}
=\frac{D_n^3(\bm x;\bm q)}{D_n^2(\bm x)}
\in[q_{\min},q_{\max}].
\end{equation}
Given a completed M3 solution, let
$\widehat H_n=D_n^2(\bm x^{\Mthree})$ and re-solve M2 with
$\widehat{\bm H}$. Then
$\mathcal V^{\MtwoMatched}\lexleq\mathcal V^{\Mthree}$ and
\begin{equation}
S^{\Mtwo}-S^{\Mthree}
=\left(S^{\Mtwo}-S^{\MtwoMatched}\right)
-\left(S^{\Mthree}-S^{\MtwoMatched}\right),
\label{eq:report_matched_decomp}
\end{equation}
where the second bracket is nonnegative.
\end{theorem}

\begin{IEEEproof}
For $z_{n,t}=x_{n,t}/P_n\geq0$,
$q_{\min}\sum_t\Delta_tz_{n,t}
\leq\sum_t\Delta_tq_{n,t}z_{n,t}
\leq q_{\max}\sum_t\Delta_tz_{n,t}$, proving the set sandwich. The used rate
is a convex combination of rates with weights proportional to
$\Delta_tz_{n,t}$. The completed M3 trajectory is feasible for M2-MATCHED by
the definition of $\widehat H_n$, which proves the value ordering. Adding and
subtracting $S^{\MtwoMatched}$ gives the decomposition.
\end{IEEEproof}

The first bracket in \eqref{eq:report_matched_decomp} is the exposure or
allowance effect and need not be positive. The second is the nonnegative shape
penalty. The matched-efficiency ratio
\begin{equation}
\eta_{\mathrm{match}}
=100\frac{S^{\Mtwo}-S^{\Mthree}}
{S^{\Mtwo}-S^{\MtwoMatched}}
\end{equation}
is reported only when the denominator exceeds the frozen numerical threshold.
M2-MATCHED is ex post and not deployable.

\subsection{Control cards}

The control family is fixed before held-out evaluation. To isolate state
ordering from arbitrary card amplitude, every nonselected ordering control uses
the selected rate vector as its amplitude reference. Given control-specific
fitting buckets with exposure weights $W_{g,b}^{\mathrm{ctrl}}$, its rates solve
\begin{align}
\min_{\bm q_g}\quad
&\sum_{b=1}^3W_{g,b}^{\mathrm{ctrl}}
(q_{g,b}-q_{g,b}^{\star})^2
\label{eq:report_control_projection}\\
\mathrm{s.t.}\quad
&q_{g,1}\geq q_{g,2}\geq q_{g,3},\quad
q_{\min}\leq q_{g,b}\leq q_{\max},
\notag\\
&\sum_bW_{g,b}^{\mathrm{ctrl}}q_{g,b}
=\sum_bW_{g,b}^{\mathrm{ctrl}}.
\notag
\end{align}
If fewer than two control buckets have positive exposure, that direction is
flat. This projection changes the selected amplitude only as required by the
common bounds, order, and reference normalization.

\begin{enumerate}[leftmargin=*,itemsep=2pt]
\item Flat M2: $q=1$.
\item PTDF net-relief score: for a one-MW admissible probe
$\delta_{n,t}=\min\{1\ \mathrm{MW},\bar x_{n,t}\}$, define the directed-overload
functional
\begin{equation}
\Phi_t(\bm f)=\sum_{c,r}
\left(\pos{f_{r,t}^c-\bar F_r^c}
+\pos{-f_{r,t}^c-\bar F_r^c}\right),
\end{equation}
and score
\begin{equation}
s_{n,t}^{\mathrm{PTDF}}
=\frac{\Phi_t(\bm f^{\mathrm{req}})
-\Phi_t(\bm f^{\mathrm{req}}+\delta_{n,t}\bm G_{:,n})}
{\delta_{n,t}}.
\label{eq:report_ptdf_score}
\end{equation}
Gross relieved overload, induced or worsened overload, and net relief are
reported separately. The score is a transparent control and never a learned
feature.
\item Inverse ordering: replace the selected score by its negative, rebuild the
fitting tertiles within direction, and apply
\eqref{eq:report_control_projection}.
\item Shuffle distribution: 32 independent within-direction permutations of the
frozen selected-score ordering. The exact seeds are
$118001,118002,\ldots,118032$. Each permutation is generated on stable
participant--time IDs, frozen before held-out access, and mapped through
\eqref{eq:report_control_projection}. Each shuffle reports $S$, $E$, physical
FDE, and burden.
\end{enumerate}

\section{Dutch Data, Chronological States, and Physical Case}
\label{sec:report_dutch_data}

\subsection{Network sources and layer separation}

The physical architecture is based on the open Dutch high-voltage models of
Zomerdijk et al.\ \cite{Zomerdijk2022,ZomerdijkData}. The detailed 383-bus
mixed-voltage model and the reduced 45-bus 220/380-kV model are extracted from
the same versioned source release but are assigned different scientific roles.
The detailed model provides physical station and device locations, participant
localization, detailed sensitivities, and the 2018 reference benchmark. The
reduced model is the only network used by the annual and rolling entitlement
optimization.

Each source extraction records: original file and release identifier, SHA-256,
software version, bus and element IDs, voltage levels, in-service flags,
parallel identifiers, line parameters, transformer parameters, tap positions,
phase shifts, ratings, shunts, generators, loads, external grids, and any
unsupported element. No downstream script is permitted to reconstruct the
network from a different source path.

\subsection{Component data and provenance}

For every province $p$ and time $t$, the state builder retains
\begin{equation}
D_{p,t},\qquad PV_{p,t},\qquad W_{p,t}^{\mathrm{on}},\qquad
W_{p,t}^{\mathrm{off}}.
\end{equation}
Each component records units, native resolution, timezone, source, revision
policy, interpolation or gap treatment, annual scaling, and scenario/forecast/
realization status. The builder does not create one residual-load trajectory and
redistribute it as a nodal injection.

The frozen V3 source registry uses the logical inputs in
Table~\ref{tab:report_source_registry}. Machine-specific absolute paths are kept
in a local companion file and are excluded from the scientific hash. The
published registry stores the logical role, file name, byte count, SHA-256, and
provenance class.

\begin{longtable}{@{}p{0.18\textwidth}p{0.35\textwidth}p{0.41\textwidth}@{}}
\caption{Logical V3 source registry. Exact hashes are carried by
\texttt{SOURCE\_REGISTRY.json}.}
\label{tab:report_source_registry}\\
\toprule
Role & Frozen logical file & Scientific status \\
\midrule
\endfirsthead
\toprule
Role & Frozen logical file & Scientific status \\
\midrule
\endhead
Reduced policy network & \path{AGGREGATED_EHV/network.current.p} & Zomerdijk-derived 45-bus candidate \\
Detailed network & \path{FULL_383/network.current.p} & localization, device inventory, reduction checks \\
Ratings & \path{branch_rating_ensemble.csv} & declared central line/transformer ratings \\
Component panel & \path{provincial_residual_load_2025_2035_learned_optionB.csv} & controlled scenario; not a D--1 forecast \\
Bus weights & \path{component_prior_weights.csv} & technology-specific province/landing-point weights \\
Reference allocations & \path{component_allocations.parquet} & public-state renewable allocation reference \\
Reference dispatch & \path{dispatch.parquet} & public-data device dispatch reference \\
Participant crosswalk & \path{applicant_direction_crosswalk_central.csv} & pre-outcome 114-position crosswalk \\
Participant contracts & \path{participants.parquet} & capacity, direction, technology, readiness, location \\
Requests & \path{participant_requests_hourly_YYYY.parquet} & controlled hourly scenario, 2025--2031 \\
2018 states/flows & \path{dutch2018_hidden_states.parquet}; \path{dutch2018_hidden_branch_flows.parquet} & independent mapped-row DC reference \\
\bottomrule
\end{longtable}

The component panel passes only when timestamps exactly match the frozen grid,
all values are finite with the declared signs, provincial sums reproduce the
national totals within tolerance, every interpolation and scaling operation is
listed, and no forecast-vintage claim is made for a scenario or realization
series.

\subsection{Technology-specific bus allocation}

Let $\bm W^D$, $\bm W^{PV}$, $\bm W^{\mathrm{on}}$, and
$\bm W^{\mathrm{off}}$ map source regions or landing points to retained buses.
Every matrix is nonnegative and each source column sums to one. Demand weights
are based on the documented load-location proxy; PV and onshore-wind weights
use technology-specific installed-capacity or generator-location evidence;
offshore wind is mapped first to documented landing points and then to the
associated EHV connection buses. A missing landing-point mapping is a hard
provenance failure, not a reason to spread offshore wind uniformly.

The nondispatchable active injection is
\begin{equation}
\bm p_t^{\mathrm{ND}}
=-\bm W^D\bm D_t
+\bm W^{PV}\bm{PV}_t
+\bm W^{\mathrm{on}}\bm W_t^{\mathrm{on}}
+\bm W^{\mathrm{off}}\bm W_t^{\mathrm{off}}.
\label{eq:report_pND}
\end{equation}
The allocation audit checks column sums, bus IDs, technology totals, national
conservation, and exact agreement between component and nodal sums.

\subsection{Dispatch and boundary balancing}

Let $\bm C$ map dispatchable generators and net-boundary devices to buses,
$\bm g_t$ their active injections, and
$[\underline{\bm g}_t,\overline{\bm g}_t]$ their available ranges. The
reference vector $\widetilde{\bm g}_t$ and positive diagonal weight matrix
$\bm W_g$ are frozen from the public scenario and fitting-year device data.
The background builder first solves the common lexicographic problem
\begin{align}
\min_{\bm g_t,\bm s_t^{\mathrm{bg}}}\quad
&S_t^{\mathrm{bg}}
:=\sum_{c,r}(s_{r,t,c}^{\mathrm{bg},+}
+s_{r,t,c}^{\mathrm{bg},-}),
\label{eq:report_bg_stage1}\\
\mathrm{s.t.}\quad
&\underline{\bm g}_t\leq\bm g_t\leq\overline{\bm g}_t,
\notag\\
&\one^\top(\bm p_t^{\mathrm{ND}}+\bm C\bm g_t)=0,
\notag\\
&\text{the retained affine thermal rows for the background state},
\notag
\end{align}
then, on the minimum-$S_t^{\mathrm{bg}}$ face, minimizes
\begin{equation}
\frac12(\bm g_t-\widetilde{\bm g}_t)^\top
\bm W_g(\bm g_t-\widetilde{\bm g}_t).
\label{eq:report_bg_stage2}
\end{equation}
The frozen background is
\begin{equation}
\bm p_t^{\mathrm{bg}}
=\bm p_t^{\mathrm{ND}}+\bm C\bm g_t^{\star}.
\end{equation}
Thus the builder makes the background as secure as the declared public device
set permits without altering component totals, and it remains independent of
M1--M3.

External grids are not modeled as mutually independent free injections. Each
synchronous boundary group has one signed net-interchange quantity and a
frozen nonnegative spatial distribution vector. The formulation therefore
cannot increase import at one boundary and export at another merely to relieve
an internal line while preserving the same national net balance. When the
source scenario prescribes a boundary sign or interchange total, that
constraint is retained.

The background artifact reports $S_t^{\mathrm{bg}}$, device bounds and
activity, net interchange, nodal balance residual, and row violations. If
$S_t^{\mathrm{bg}}>0$, the residual is not hidden; the requested-state and
post-policy results report both the total and the increment relative to the
background.

\subsection{Reactive-power states for engineering reconstruction}

The parent entitlement model is active-power DC and does not use reactive power.
A separate immutable $\bm q_t^{\mathrm{bg}}$ artifact is nevertheless required
for the AC engineering levels. Existing source-model load Q, shunts, generator
voltage set points, Q bounds, transformer taps, and external-grid controls are
preserved whenever available. For a load lacking a complete Q chronology, its
reactive demand is reconstructed from a frozen fitting-period or source-model
$Q/P$ ratio at the same device or documented class; a missing ratio is a
provenance failure and is not replaced silently by zero. Controllable generators
remain PV devices with their source voltage targets and Q limits, whereas
noncontrollable PQ devices retain their documented power factor. The AC source
manifest records every device-level rule, imputation, limit, and hash.

Prospective participant P/Q changes are not embedded in the parent action
matrix. In selected-state AC tests they are added under project-specific
reactive data when available and otherwise under the predeclared unity,
$0.98$ lagging, and $0.98$ leading power-factor sensitivities. This separation
prevents an assumed Q response from altering the annual M1--M3 optimization.

\subsection{Residual load as a descriptor}

Provincial residual load is
\begin{equation}
RL_{p,t}
=D_{p,t}-PV_{p,t}-W_{p,t}^{\mathrm{on}}-W_{p,t}^{\mathrm{off}}.
\end{equation}
It is used to summarize net provincial demand pressure and, with authenticated
vintages, to form M3 features. It is not a branch flow, a nodal power-flow
solution, or the input to a uniform bus allocation. The state builder stores
both the separate components and the derived descriptor so that conservation
and feature provenance can be audited independently.

\subsection{2018 Zomerdijk reference checks}

The 2018 release supplies nodal injections and mapped branch flows for an
independent reconstruction check. The intact affine evaluator is applied to the
2018 injections and compared row by row with the stored branch-flow reference.
This is an exact software/network-construction test: bus and row IDs, orientation,
ratings, balance, MW flow, sign, limiting element, and false-safe/false-overload
classifications must pass their numerical tolerances.

A separate plausibility summary compares total and voltage-level demand,
technology-level generation, branch-loading quantiles, overlap among stressed
interfaces, and available voltage/reactive ranges. Exact chronological identity
with 2025--2031 is neither expected nor claimed because the scenario years and
component paths differ. The reference checks are fixed before policy outcomes
and are not tuned to improve M1--M3 comparisons.

\section{Reduced-Network Contingencies and Participant Mapping}
\label{sec:report_network_mapping}

\subsection{Affine PTDF construction with phase shifts}

For state $c$, let $\bm B_{\mathrm{bus}}^c$ be the reduced nodal susceptance
matrix after selecting a reference bus, $\bm B_f^c$ the branch-angle matrix,
and $\bm b_{\mathrm{shift}}^c$ the fixed phase-shift contribution. The DC
equations are
\begin{align}
\bm B_{\mathrm{bus}}^c\bm\theta_t^c
&=\bm p_t-\bm d_{\mathrm{shift}}^c,
\notag\\
\bm f_t^c
&=\bm B_f^c\bm\theta_t^c+\bm b_{\mathrm{shift}}^c.
\end{align}
Eliminating $\bm\theta_t^c$ gives
\begin{equation}
\bm f_t^c=\bm H^c\bm p_t+\bm h^c.
\end{equation}
This construction follows the standard DC approximation while retaining the
affine phase-shift term \cite{Stott2009}. The implementation verifies invariance
to the chosen reference bus and stores both $\bm H^c$ and $\bm h^c$. The row
orientation follows the source
from-bus/to-bus convention; the reverse inequality is represented by sign, not
by silently changing the branch ID.

\subsection{Component and rating inventory}

The network inventory assigns each line and transformer:
physical element ID, source table and index, from/to bus, voltage classes,
parallel-circuit identifier, in-service status, reactance, tap ratio,
phase-shift angle, continuous rating, any emergency rating, and support status.
A rating is accepted only with known units and conversion. If only a single
rating exists, it is used consistently and described as the declared public
rating. No emergency factor is invented.

\subsection{Contingency enumeration}

Candidate states consist of the intact topology and one-outage states generated
from supported lines and transformers. A candidate is retained if:
\begin{enumerate}[leftmargin=*,itemsep=2pt]
\item the source element represents one identifiable modeled physical element;
\item removing it leaves the modeled EHV network connected and the DC matrix of
the retained island nonsingular after reference selection;
\item every remaining monitored element has a valid orientation and rating;
\item the affine evaluator and independent DC reconstruction pass their
structural tests.
\end{enumerate}
If parallel physical circuits are listed separately, each is a separate
candidate. If the source contains one equivalent branch for a multi-circuit
corridor, removing the equivalent is not called a single-circuit N--1 event;
it is excluded or labeled an equivalent-corridor stress state. Islanding,
unsupported devices, missing ratings, and unresolved semantics are listed in
the exclusion registry.

The retained claim is therefore limited to ``retained non-islanding
single-element contingencies represented by the reduced EHV model.'' The
report publishes the complete registry rather than summarizing it only by a
count.

\subsection{Direct participant transactions}

For a direct position $n$ connected at retained bus $b(n)$, the participant
side of the transaction is one-hot. The counter-injection side uses the frozen
upward or downward balancing key consistent with the direction of the
prospective request. The resulting column $\bm a_n$ is balanced exactly and
retains the same balancing-key hash used in the detailed and reduced
sensitivity calculations.

The detailed model uses the corresponding physical station. The direct
comparison is not assumed to be exact because network reduction can alter
sensitivities; it provides the empirical reference distribution against which
the multibus embeddings are judged.

\subsection{Multibus embedding calibration}

For a position whose detailed connection point is absent from the reduced
model, let $\mathcal B_n^{\mathrm{sup}}$ be the predeclared nearby or boundary
support buses and $\bm w_n$ their nonnegative embedding weights. They satisfy
\begin{equation}
\bm w_n\geq0,
\qquad
\one^\top\bm w_n=1.
\label{eq:report_embedding_simplex}
\end{equation}
Let $\bm g_{n}^{383,c}$ be the detailed participant effect on the common
monitored EHV interface map, normalized by the corresponding rating. Let
$\bm G_{\mathcal B_n}^{45,c}$ collect the reduced effects of the support-bus
transactions under the identical balancing convention. The calibration is
\begin{align}
\min_{\bm w_n}\quad
&\sum_{c\in\mathcal C^{\mathrm{map}}}
\norm{\bm\Omega_c^{1/2}
(\bm G_{\mathcal B_n}^{45,c}\bm w_n-\bm g_n^{383,c})}_2^2
+\lambda_w\norm{\bm w_n-\bm w_n^0}_2^2
\label{eq:report_embedding_qp}\\
\mathrm{s.t.}\quad
&\eqref{eq:report_embedding_simplex},
\notag
\end{align}
where $\bm\Omega_c$ is a frozen row-weight matrix, $\bm w_n^0$ is a
geographical/electrical prior, and $\lambda_w$ is fixed before policy
outcomes. Only network structure and fitting-period device data may enter this
calibration.

\subsection{Detailed-to-reduced validation metrics}

For each position and common state, let
$\bm z_n^{D,c}$ and $\bm z_n^{R,c}$ be the detailed and reduced
rating-normalized sensitivity vectors on exactly aligned row IDs. Define
\begin{align}
e_{2,n}^c
&=\frac{\norm{\bm z_n^{R,c}-\bm z_n^{D,c}}_2}
{\max\{\norm{\bm z_n^{D,c}}_2,10^{-9}\}},
\notag\\
e_{\infty,n}^c
&=\frac{\norm{\bm z_n^{R,c}-\bm z_n^{D,c}}_\infty}
{\max\{\norm{\bm z_n^{D,c}}_\infty,10^{-9}\}}.
\label{eq:report_embedding_errors}
\end{align}
Sign agreement is evaluated on material rows satisfying
$|z_{n,r}^{D,c}|\geq0.05\norm{\bm z_n^{D,c}}_\infty$. Top-$k$ overlap uses
$k=\min\{10,|\mathcal R_c^{\mathrm{map}}|\}$ and the largest absolute
sensitivities.

The 79 directly mapped positions provide the empirical benchmark for the
unavoidable detailed-to-reduced sensitivity error of the 45-bus
representation. Before any policy outcome is generated, the direct-position
sample defines reference thresholds for each reduction metric. Let
$\tau_2$, $\tau_\infty$, $\tau_{\mathrm{sign}}$, and $\tau_{\mathrm{top}}$
denote respectively the frozen upper-tail error thresholds and lower-tail
agreement thresholds obtained from the direct-position benchmark. A multibus
embedding is admitted to the primary participant set only if
\begin{equation}
\max_c e_{2,n}^{c}\leq\tau_2,\qquad
\max_c e_{\infty,n}^{c}\leq\tau_\infty,
\qquad
\min_c\mathrm{SignAgree}_{n}^{c}\geq\tau_{\mathrm{sign}},
\qquad
\min_c\mathrm{TopK}_{n}^{c}\geq\tau_{\mathrm{top}}.
\label{eq:report_embedding_gate}
\end{equation}
The thresholds are computed exclusively from the direct-mapping reduction
benchmark and frozen before any M1--M3 policy outcome is generated. They
therefore test whether an electrical multibus embedding introduces materially
greater reduction error than that already present for an exactly localized
position in the reduced model; they are not presented as universal
network-reduction standards.

If all 35 multibus positions pass this gate, all 114 positions form the
primary participant set. Otherwise the 79 direct positions form the
fail-closed primary case, while the complete 114-position case is reported
only as a sensitivity. Partial post-outcome inclusion of multibus positions
is prohibited.

The frozen benchmark uses the 95th percentile of the direct-position
worst-state $e_2$ and $e_\infty$ errors as $\tau_2$ and $\tau_\infty$, and the
5th percentile of the corresponding sign-agreement and top-$k$ overlap
statistics as $\tau_{\mathrm{sign}}$ and $\tau_{\mathrm{top}}$.
\subsection{Primary physical-case freeze}

The physical case can be frozen only after the following artifacts pass:
network source inventory; device and rating audit; phase-shift audit;
component-data and allocation audit; background balance and security audit;
contingency registry; participant-action balance audit; exact
$\bm G^c=\bm H^c\bm A$ identity; detailed-to-reduced validation and primary-set
decision; and time-grid/request/mask alignment. The freeze manifest records
every input path, byte size, hash, identifier order, unit, and software version.
No experiment runner may rebuild or mutate a physical-case object.

\section{Annual, Rolling, and Control Experiment Protocol}
\label{sec:report_experiments}

\subsection{Chronological split and information lock}

The chronological split is fixed as follows:
\begin{equation}
\underbrace{2025\text{--}2027}_{\text{fit}}
\quad\longrightarrow\quad
\underbrace{2028}_{\text{operational validation}}
\quad\longrightarrow\quad
\underbrace{2029\text{--}2030}_{\text{held-out evaluation}}
\quad\longrightarrow\quad
\underbrace{2031}_{\text{separate confirmation}}.
\label{eq:report_split}
\end{equation}
No 2029--2031 outcome may select or alter the feature family, predictor,
preprocessing, score thresholds, card rates, rate bounds, card amplitude,
regularization, participant set, contingency set, shuffle seeds, physical
case, acceptance thresholds, or manuscript claim rule. The held-year access
audit is executed before card serialization and again before evaluation.

\subsection{Allowance regimes and annual cases}

Two regimes have distinct scientific roles.

\begin{table}[t]
\centering
\caption{Frozen annual experiment contract.}
\label{tab:annual_contract}
\begin{tabularx}{\textwidth}{@{}p{0.10\textwidth}p{0.31\textwidth}X@{}}
\toprule
Nominal share & Required annual cases & Scientific role \\
\midrule
$\rho=0.15$
& M1-LP, M2, allowance-specific frozen M3
& Primary duration-to-depth comparison and strict M1 certificate. Rolling M1
supplies the feasible annual upper endpoint. If M3 is flat, it is reported as
M2. \\
$\rho=0.03$
& M2, allowance-specific frozen M3, M2-MATCHED, PTDF/inverse/32-shuffle controls
& State-conditioning mechanism study with physical exposure and shape separated.
\\
\bottomrule
\end{tabularx}
\end{table}

Each allowance regime has its own teacher, validation, and card freeze. Every
same-scale case uses the identical physical-case and allowance-vector hashes.
M2-MATCHED is the only exception and records the participant-specific vector
generated by the completed M3 trajectory. A different participant set,
contingency set, background state, rating, time grid, or allowance vector defines
a sensitivity rather than the same-scale theorem comparison.

\subsection{Rolling controller}

The primary controller uses a 48-h look-ahead and implements the first 24 h.
At epoch $\tau$, it solves the Stage-I/II sequence over
$\mathcal L_\tau$, implements $\mathcal I_\tau$, and updates
\begin{equation}
B_{n,\tau+1}^M
=B_{n,\tau}^M-
\sum_{t\in\mathcal I_\tau}\Delta_t d_{n,t}^M,
\end{equation}
where $d^1=u$, $d^2=x/P$, and $d^3=qx/P$. All policies have identical
look-ahead, implementation block, forecasts, requests, physical rows, and
solver tolerances. The primary controller has no terminal reserve. A paced or
terminal-value design, if studied, is clearly labeled as a separate
sensitivity and cannot replace the primary result.

The rolling-M1 epoch is accepted only when its implemented MILP action is
feasible, its debit is replayed, its balance remains nonnegative, and its
checkpoint contains the exact prior balance and source hashes. Concatenation
checks timestamp coverage, nonoverlap, action continuity, and the annual scalar
upper-bound conditions.

\subsection{Common metrics}

For every policy and year, the following are reported on the same units and
population:
\begin{enumerate}[leftmargin=*,itemsep=2pt]
\item $S$ and $E$;
\item materially touched duration under the M1 materiality definition;
\item physical FDE $\sum_{n,t}\Delta_tx_{n,t}/P_n$;
\item weighted debit $\sum_{n,t}\Delta_tq_{n,t}x_{n,t}/P_n$;
\item participant-level and aggregate restricted MWh;
\item ex-post price-weighted restriction burden
$\sum_{n,t}\Delta_t\pi_tx_{n,t}$, reported only as a diagnostic;
\item restriction-episode count, isolated one-interval episodes, and maximum
episode duration;
\item affected participants and binding entitlement accounts;
\item maximum overload in MW and loading fraction;
\item insecure intervals and violated row/state/interval tuples;
\item background, requested-state, and post-policy residual indices;
\item rating-normalized residual severity
\begin{equation}
S^{\mathrm{norm}}
=\sum_t\Delta_t\sum_{c,r}
\frac{s_{r,t,c}^{+}+s_{r,t,c}^{-}}{\bar F_r^c};
\end{equation}
\item runtime, model dimensions, retained rows, closure rounds, solver status,
bounds, and replay residuals.
\end{enumerate}
The theorem remains stated in terms of $S$. Normalized severity is a physical
sensitivity diagnostic and never replaces the objective after outcomes are
seen.

\subsection{Source-of-separation decomposition}

The M1-LP--M2 debit difference
\begin{equation}
\sum_{n,t}\Delta_tx_{n,t}
\left(\frac1{\bar x_{n,t}}-\frac1{P_n}\right)
\end{equation}
is decomposed by bins of request headroom
$\bar x_{n,t}/P_n$. This identifies whether the set expansion is concentrated
in low-request coordinates, near-full-request coordinates, particular
participant directions, or particular periods. The decomposition is an exact
accounting identity evaluated on a fixed trajectory; it is not a causal
attribution.

\subsection{M3 control interpretation}

Each M3 or control card is reported jointly with its physical FDE and
M2-MATCHED result. The primary mechanism quantities are
\begin{align}
\Delta S^{\mathrm{nom}}
&=S^{\Mtwo}-S^{\Mthree},
\notag\\
\Delta S^{\mathrm{allow}}
&=S^{\Mtwo}-S^{\MtwoMatched},
\notag\\
P^{\mathrm{shape}}
&=S^{\Mthree}-S^{\MtwoMatched}\geq0.
\end{align}
A lower nominal $S$ with materially greater physical FDE is described as an
entitlement/exposure deformation unless matched exposure shows a remaining
allocation gain. Shuffle controls are summarized as a distribution; no single
arbitrary shuffle is treated as the random baseline.

\section{Engineering Validation Protocol}
\label{sec:report_engineering}

\subsection{Selection manifest}

Engineering cases are selected from frozen parent-model outputs before any
independent DC or AC outcome is viewed. For each evaluation year, the manifest
selects two hours from each of three categories:
\begin{enumerate}[leftmargin=*,itemsep=2pt]
\item largest unrestricted-request residual stress;
\item largest rolling M1--M2 flow difference;
\item largest annual M2--M3 flow difference.
\end{enumerate}
Ties are resolved by timestamp and stable policy ID, duplicates are removed, and
a deterministic fill rule preserves the target count where possible. For each
hour, the registry includes the intact state, the worst requested-state outage,
and the state with the largest cross-policy flow spread. The manifest stores the
balanced injection and selection-source hashes and is reused unchanged by all
engineering levels.

\subsection{Level 1: independent DC reconstruction}

The first validation evaluates the exact same balanced injection by two
independent paths:
\begin{enumerate}[leftmargin=*,itemsep=2pt]
\item the frozen affine evaluator $\bm H^c\bm p+\bm h^c$;
\item a separately coded DC power flow assembled from the network element data,
reference-bus convention, tap ratios, and phase shifts.
\end{enumerate}
A permanent mapping aligns network state, element, orientation, and rating.
The audit reports:
\begin{align*}
&\text{MAE and RMSE in MW},\quad
\text{maximum absolute error},\quad
\text{rating-normalized errors},\\
&\text{sign agreement},\quad
\text{limiting-element agreement},\quad
\text{top-5 and top-10 overlap},\\
&\text{false-safe and false-overload classifications for all rows and rows
within 10 percentage points of rating}.
\end{align*}

Because both paths implement the same DC model, the acceptance gate is
numerical rather than statistical: exact row/contingency IDs, maximum error no
larger than
$\max\{10^{-4}\ \mathrm{MW},10^{-8}\bar F_r^c\}$, 100\% sign agreement after
excluding flows below $10^{-6}$ MW, identical limiting element outside ties,
and zero false-safe/false-overload classifications beyond tolerance. Failure
invalidates row-level DC policy claims until the mapping or evaluator is
corrected.

\subsection{Level 2: selected-state nonlinear consistency and stress}

The selected-state AC audit uses the same network release and mapped topology
as the detailed public model. It preserves all available:
PV and slack controls, generator voltage targets, reactive-power bounds,
shunts, transformer tap positions and phase shifts, in-service statuses, and
ratings. The active-power injection is the frozen policy state. Participant
reactive-power changes are evaluated under unity power factor and 0.98 lagging
and leading sensitivities unless project-specific reactive data are available.

The numerical attempt sequence changes only initialization and declared
control handling; it does not change the active-power injection to force
convergence. It proceeds through: (i) exact reconstructed PQ injections with the
source slack treatment; (ii) the same active injections with source-supported
PV voltage control and Q limits restored; and (iii) a stable flat/DC-angle
initialization under the same device equations. Standard PV-to-PQ switching on
Q-limit hits is recorded rather than hidden. No attempt scales load, generation,
or the parent policy action. Every attempt and failure reason is stored. The
implementation uses a versioned AC solver such as pandapower
\cite{Thurner2018}. The report separates:
\begin{enumerate}[leftmargin=*,itemsep=2pt]
\item \emph{convergence}: whether the nonlinear equations solve;
\item \emph{voltage feasibility}: whether declared or sensitivity voltage bands
are respected;
\item \emph{thermal feasibility}: whether AC MVA loading respects its AC rating;
\item \emph{reactive feasibility}: whether generator and device Q bounds are
respected.
\end{enumerate}
The bands $[0.95,1.05]$ p.u. and $[0.90,1.10]$ p.u. may be reported as narrow
and restoration sensitivities, respectively, but are not called statutory
limits unless the source establishes that status.

Mapped AC active power $P^{\mathrm{AC}}$ is compared with DC active power
$P^{\mathrm{DC}}$ in MW on the same branch and orientation. The report gives
MAE, RMSE, normalized MAE, identity-prediction $R^2$, Spearman rank, sign
agreement, near-limit classification agreement, and limiting-element
agreement. DC MW is never compared on a one-to-one line with AC MVA.

For a selected case, the ordering statistic is total normalized thermal
exceedance,
\begin{align}
J^{\mathrm{DC}}
&=\sum_e\pos{\abs{P_e^{\mathrm{DC}}}/\bar F_e^{\mathrm{DC}}-1},
\notag\\
J^{\mathrm{AC}}
&=\sum_e\pos{\abs{S_e^{\mathrm{AC}}}/\bar S_e^{\mathrm{AC}}-1}.
\label{eq:report_ac_ordering}
\end{align}
The two quantities are used only for within-model policy ordering; they are not
placed on a MW--MVA one-to-one line. Voltage and reactive violations remain
separate outcomes. The frozen ordering-tie tolerance is
$\tau_J=10^{-4}$ in normalized-exceedance units. For each predeclared policy
pair $(A,B)$, a reversal occurs when both within-model differences exceed
$\tau_J$ in magnitude and the sign of
$J_A^{\mathrm{DC}}-J_B^{\mathrm{DC}}$ differs from that of
$J_A^{\mathrm{AC}}-J_B^{\mathrm{AC}}$. A DC-nontied pair that becomes AC-tied
is reported separately as loss of discrimination, not as preservation. The
wording ``nonlinear detail does not reverse the selected-state ordering'' is
permitted only when every DC-nontied pair either retains its sign or is reported
explicitly as an AC tie, with no reversal. Otherwise, the report gives the
reversal count and uses restricted wording about broad association or mixed
consistency.

\subsection{Level 3: optional all-hour background AC reconstruction}

The optional all-hour experiment reconstructs the intact background P/Q state
for every hour of the configured year. It validates the chronology and
physical plausibility of $\bm p_t^{\mathrm{bg}}$, not the M1--M3 actions. It is
checkpointed in seven-day partitions and reports convergence, voltage
statistics, reactive-limit activity, active-flow AC/DC metrics, and hours with
violations. It does not establish annual policy AC feasibility, AC-OPF
optimality, or contingency security.

\subsection{Common redispatch sensitivity audit}

For a selected frozen policy state, let $\bm R\bm w$ be a balanced remedial
set with $\one^\top\bm R=0$ and common bounds
$\underline{\bm w}\leq\bm w\leq\overline{\bm w}$. The audit first minimizes
$S$ with $x$ fixed to the parent policy, then minimizes the common
redispatch effort
\begin{equation}
C^{\mathrm{RD}}(\bm w)=\sum_j c_j|w_j|
\end{equation}
on the minimum-$S$ face. The same action matrix, costs, bounds, and objective
are used under M1, M2, and M3. The primary question is whether common
substitute remedies reverse the accounting-rule ordering.

An auxiliary restoration sensitivity may allow additional restriction
$\Delta x\geq0$. Such restriction remains subject to the original action bound
and is debited against the remaining entitlement using the applicable rule.
A result that ignores this additional debit is not admissible evidence.
Redispatch claims require complete source, input, output, and status artifacts;
an incomplete audit is reported only as a limitation.

\section{Post-Freeze Rate and Temporal Sensitivity}
\label{sec:report_sensitivity}

\subsection{Rate-bound sensitivity}

The primary card uses $[0.5,1.5]$. The predeclared tight and wide bounds are
$[0.7,1.3]$ and $[0.25,2.0]$. These cards use the selected frozen score model,
preprocessing, thresholds, and training rows. Only the rate QP is re-solved
under the alternate predeclared bounds, and every annual case is solved
exactly. The report always places changes in $S$ beside changes in restricted
MWh, physical FDE, burden, and M2-MATCHED.

\subsection{Card-amplitude path}

For the serialized selected vector $\bm q^\star$, define
\begin{equation}
\bm q^{\mathrm{raw}}(\gamma)
=\one+\gamma(\bm q^\star-\one).
\label{eq:report_gamma_raw}
\end{equation}
The endpoint $\gamma=0$ is the exact flat M2 card and $\gamma=1$ is the exact
serialized selected card; neither endpoint is replaced by a second numerical
projection. For a nonendpoint $\gamma$, each direction solves
\begin{align}
\min_{\bm q}\quad
&\sum_{b=1}^3W_{g,b}^{\mathrm{ref}}
(q_b-q_b^{\mathrm{raw}}(\gamma))^2
\label{eq:report_gamma_projection}\\
\mathrm{s.t.}\quad
&q_1\geq q_2\geq q_3,
\quad 0.25\leq q_b\leq2.0,
\notag\\
&\sum_bW_{g,b}^{\mathrm{ref}}q_b
=\sum_bW_{g,b}^{\mathrm{ref}}.
\notag
\end{align}
The frozen path uses
\begin{multline}
\Gamma=\{0,0.25,0.35,0.45,0.55,0.65,0.70,0.80,0.90,1.00,\\
1.10,1.20,1.30,1.40,1.50,1.60,1.70,1.80,1.90,1.95,2.00\}.
\label{eq:report_gamma_grid}
\end{multline}
Every nontrivial gamma/year performs a full M3 and M2-MATCHED solve. The
$\gamma=0$ endpoint reuses the
hash-verified exact M2 solve. Performance curves are not interpolated or
prescribed.

\subsection{Fixed-policy temporal block resampling}

Seven-day circular blocks ($168$ h) are sampled from the already optimized
hourly contribution arrays using $1000$ draws, following the standard rationale
for block resampling of dependent time series \cite{Kunsch1989,PolitisRomano1994}.
The M1 chronology diagnostic uses seed $118+y$. A gamma case uses seed
$3118+1000y+\operatorname{round}(1000\gamma)$. The same block indices are
applied to all policies in a paired comparison. The procedure preserves
within-policy temporal dependence at the chosen block scale but does not
re-optimize, resample participants, or represent an unknown annual optimum. Its
bands are labeled ``fixed-policy temporal block resampling'' and never
``optimization-gap confidence interval,'' ``generalization interval,'' or
``participant bootstrap.'' The network-coupled participants are not
independently resampled.

\section{Out-of-Sample and Failure-Mode Stress Protocol}
\label{sec:report_stress}

Stress tests are predeclared diagnostics and do not redefine the primary case.
They are reported in absolute MW, MWh, participant-hours, EUR, violated hours,
and solver seconds.

\subsection{Debit-rate and allowance stress}
The selected score model and thresholds are fixed while the rate bounds are
changed to $[0.7,1.3]$ and $[0.25,2.0]$. The amplitude path in
Section~\ref{sec:report_sensitivity} tests the full declared $\gamma$ grid. Each
case reports M3, M2-MATCHED, physical FDE, restricted MWh, price burden, and the
state-shape penalty. Very wide bounds are treated as a failure-mode test of
entitlement deformation, not as an improved product by default. Allowance stress
may add predeclared shares outside $0.03$ and $0.15$, but a card is never moved
between allowance hashes.

\subsection{Topology stress}
The primary outage family is one retained supported element at a time. Additional
$N-k$ tests, where constructed, are labeled stress topologies and do not enter
the primary N--1 wording. A stress state must preserve exact element identity,
connectivity or explicit island treatment, ratings, and an independent DC map.
Unbounded or missing thermal limits are not replaced by arbitrary large numbers;
the affected row is excluded with a recorded reason or the stress case fails.

\subsection{Objective and penalty limits}
The parent model uses exact lexicographic stages rather than a scalar penalty
weight between security and restricted energy. There is therefore no
``infinite penalty'' parameter to send to a limit. A weighted-sum sensitivity
may be used only as a separate diagnostic and may not replace the Stage-I/II
result. Slack upper bounds, if introduced for numerical tests, must be proven
inactive at the accepted solution or the case is not comparable.

\subsection{Information and distribution shift}
OOD diagnostics stratify 2029--2031 by request fraction, realized provincial
residual load, season, direction, technology, and distance from the fitting
feature support. Realized residual load is used only after the card is frozen.
The report gives absolute policy outcomes and control ranks within each stratum.
No OOD diagnostic is used to retune the card or claim a probabilistic guarantee.

\subsection{Engineering failure cases}
Independent DC mapping failures, AC nonconvergence, voltage violations, thermal
violations, Q-limit activity, and policy-order reversals are retained as results.
The analysis does not discard difficult selected states. A failed selected-state
AC case blocks only the corresponding nonlinear wording; it does not alter the
parent DC policy after the fact.

\section{Performance-Preserving Implementation Contract}
\label{sec:report_performance}

Performance engineering may change representation and solver settings but may
not change the mathematical experiment. The current V3 model exploits active-
hour compression, positive participant--hour coordinate compression, a low-
dimensional factorization of participant flow effects, sparse matrix transfer,
Stage-I to Stage-II basis transfer, and rolling overlap starts. The exact row-
generation theorem in Section~\ref{sec:report_certification} permits a small
initial row seed while retaining the full row universe in closure and replay.

The performance baseline is the validation year 2028. Before held-out reruns,
the reference and accelerated M1-LP and M2 implementations are compared on:
$S$, $E$, allowance and debit residuals, action bounds, complete row replay,
closed-row IDs, active hours, retained action coordinates, row counts, nonzeros,
closure rounds, build seconds, Stage-I/II seconds, replay seconds, and peak
memory. The accelerated path is accepted only if the scientific values and all
certificates agree within the frozen tolerances.

Presolve is not assumed to be beneficial. Earlier large-model tests showed that
aggressive internal presolve can duplicate sparse storage and exceed the host
memory even when it reduces the final row count. Candidate profiles therefore
include no-presolve dual simplex, automatic method, and selected presolve/thread
settings, each measured on the same 2028 model. A solver parameter is frozen only after both time and memory are recorded. The qualified primary profile uses $10^{-8}$ feasibility/optimality tolerances. A replay-triggered recovery pass may increase numerical care and tighten those tolerances, but it does not change the declared mathematical experiment and is recorded in the stage metadata. Stage III is excluded from the primary performance comparison because it is not part of $(S,E)$.

Independent policy years may be executed concurrently only after one-case peak
memory is known and the license permits it. Concurrency changes wall-clock time,
not the experiment identity. Every case remains atomic and restart-safe.

\section{Reproducibility Architecture and Artifact Contract}
\label{sec:report_reproducibility}

\subsection{Clean V3 module architecture}

The authoritative implementation is organized by scientific responsibility:
\begin{verbatim}
revision_v3/
  tpwrs_v3/
    experiment_contract.py
    time_grid.py
    physical_case.py
    io_v3.py
    model_v3.py
    linear_model.py
    screening_v3.py
    replay_v3.py
    rolling_v3.py
    m3_v3.py
    features_v3.py
    predictors_v3.py
    card_selection_v3.py
    controls_v3.py
    gamma_v3.py
    engineering_audits.py
    metrics.py
    certification.py
    dutch/
      component_data.py
      state_builder.py
      network_reduction.py
      participant_effects.py
      contingency_builder.py
      zomerdijk_extract.py
  scripts/
    build_real_physical_case.py
    build_f0_features.py
    fit_freeze_cards.py
    run_final_m1_m2_campaign.py
    run_final_m3_core.py
    run_m3_controls.py
    run_m3_rate_bound_sensitivity.py
    run_m3_gamma_path.py
    build_selected_state_registry.py
    run_independent_dc_selected_states.py
    run_selected_ac_validation_v3.py
    aggregate_final_evidence.py
\end{verbatim}
Scientific parameters are passed through typed immutable objects. Configuration
modules may define defaults, but the solver does not infer a share, time step,
card, or result path from an untracked global.

\subsection{Required physical-case artifact tree}

A frozen physical case contains at least:
\begin{verbatim}
physical_case/<physical_case_id>/
  PHYSICAL_CASE.json
  time_grid.parquet
  buses.csv
  elements.csv
  devices.csv
  ratings.csv
  component_panel.parquet
  allocation_W_D.npz
  allocation_W_PV.npz
  allocation_W_ON.npz
  allocation_W_OFF.npz
  background_p.parquet
  background_q.parquet
  requests.parquet
  contract_duration.parquet
  restriction_eligibility.parquet
  participants.csv
  participant_A.npz
  contingency_registry.csv
  affine_H.npz
  affine_h.npz
  participant_G.npz
  source_manifest.csv
  audits/
\end{verbatim}
Every matrix has explicit row and column IDs in its metadata; sparse and dense
representations are mathematically interchangeable only when their ID maps and
hashes agree.

\subsection{Required solution artifact tree}

Each experiment case writes atomically to its own directory:
\begin{verbatim}
results/<case_id>/
  EXPERIMENT_CONTRACT.json
  STATUS.json
  MODEL_BUILD_MANIFEST.json
  STAGE1_CERTIFICATE.json
  STAGE2_CERTIFICATE.json
  SELECTOR_AUDIT.json          # when Stage III is used
  ROW_CLOSURE_LOG.csv
  solution_x.parquet
  solution_u.parquet           # M1 only
  solution_q.parquet           # M3/control cards
  participant_metrics.parquet
  hourly_metrics.parquet
  REPLAY_STATUS.json
  source_manifest.csv
  solver/
\end{verbatim}
M2-MATCHED additionally writes the actual allowance vector and its M3 source
hash. Rolling cases write epoch contracts, checkpoints, implemented blocks,
and a complete ledger.

\subsection{Software environment and deterministic behavior}

The environment capture records operating system, Python version, package
versions, solver version and license type, CPU and thread settings, random
seeds, BLAS information, timezone database, and source hashes. Continuous
teacher and card-sensitive solves use one deterministic solver thread. M1
rolling profiles may use multiple threads only when the exact profile is
recorded; reproducibility is assessed by objective, feasibility, and selected
policy tolerances rather than byte identity of solver logs.

Every write is atomic. A completed status lists all produced artifacts, byte
sizes, and SHA-256 values. A case is reusable only when its contract, code,
physical case, card, and output hashes match. Adding unrelated post-freeze
files does not invalidate tracked frozen files, but changing or removing a
tracked source is a hard provenance failure.

\subsection{Methodology-to-implementation contract}

The normative machine-readable crosswalk is distributed as
\path{TPWRS_Methodology_Implementation_Contract.csv}. It maps every
scientific object to its equation, authoritative module/function, required
input artifact, and required audit output. No legacy function is authoritative
merely because it produced an earlier result. During implementation, a missing
input is reported as a methodology contract conflict with the failed
assumption, evidence, smallest valid correction, theorem consequence, and
experimental consequence.

\section{Machine-Readable Parameter and Topology Appendices}
\label{sec:report_raw_tables}

The exact topology and parameter values are distributed as machine-readable
files rather than reconstructed from rounded prose tables. The submission data
bundle must contain, at minimum:
\begin{verbatim}
buses.csv
network_elements.csv
devices.csv
ratings.csv
contingency_registry.csv
participants.csv
component_panel.parquet
allocation_W_D.npz
allocation_W_PV.npz
allocation_W_ON.npz
allocation_W_OFF.npz
background_p.parquet
background_q.parquet
participant_A.npz
affine_H.npz
affine_h.npz
participant_G.npz
source_manifest.csv
\end{verbatim}
Each file carries explicit row and column IDs, units, source hashes, and the
physical-case hash. The bus table includes bus ID, nominal voltage, source index,
and retained status. The element table includes element ID, type, from/to bus,
parallel ID, reactance, tap ratio, phase shift, in-service status, and rating
source. The device table includes bus, technology, lower/upper active limits,
reference dispatch, boundary group, reactive data, and source rule. The
contingency registry gives one row per candidate with inclusion status and
reason. The state-matrix package carries the reduced $B_{\rm bus}$ and $B_f$
objects or the information needed to reproduce them, together with $H$ and $h$.

Generator production-cost curves are not inputs to the parent or background
objectives. The background second stage uses frozen reference-distance weights,
not market cost. No cost curve is fabricated to satisfy a generic data-table
template.

For archival submissions that require typeset tables, the following inputs are
generated directly from the frozen CSV files. Their absence does not authorize
manual transcription.
\IfFileExists{data_tables/buses_longtable.tex}{\input{data_tables/buses_longtable.tex}}{
\paragraph{Bus table.} To be generated from the frozen \texttt{buses.csv}; the
present source package did not include the row-level physical-case table.}
\IfFileExists{data_tables/elements_longtable.tex}{\input{data_tables/elements_longtable.tex}}{
\paragraph{Element table.} To be generated from the frozen
\texttt{network\_elements.csv}; the present source package did not include the
row-level topology.}
\IfFileExists{data_tables/devices_longtable.tex}{\input{data_tables/devices_longtable.tex}}{
\paragraph{Device table.} To be generated from the frozen \texttt{devices.csv}.}
\IfFileExists{data_tables/contingencies_longtable.tex}{\input{data_tables/contingencies_longtable.tex}}{
\paragraph{Contingency table.} To be generated from the frozen
\texttt{contingency\_registry.csv}.}

\section{Claim-to-Artifact Acceptance Matrix}
\label{sec:report_claims}

Table~\ref{tab:claim_matrix} fixes the evidence needed before manuscript
wording is authorized. The status evaluator is run after the final campaign;
writing is generated from its pass/fail outcomes rather than from legacy
headline values.

\begin{longtable}{@{}p{0.17\textwidth}p{0.20\textwidth}p{0.25\textwidth}p{0.30\textwidth}@{}}
\caption{Normative claim-to-artifact matrix.}
\label{tab:claim_matrix}\\
\toprule
Intended claim & Required evidence & Acceptance condition & Wording contract \\
\midrule
\endfirsthead
\toprule
Intended claim & Required evidence & Acceptance condition & Wording contract \\
\midrule
\endhead
M1-LP is the exact continuous projection
& Theorem~\ref{thm:report_hierarchy}; unit and solver regression tests
& Algebraic tests pass for zero, submaterial, full, and heterogeneous bounds
& State exact projection and qualified isolated convex hull; otherwise remove convex-hull wording. \\
M2 feasible set contains M1-LP
& Debit replay and hierarchy tests
& Same physical case and allowance hash; no positive violation of \eqref{eq:report_gap}
& State common-model inclusion; otherwise methodology contract conflict. \\
M2 or M3 is strictly better than every feasible annual M1 schedule in Stage I
& Replay-feasible Q; certified M1-LP lower bound; tolerance manifest
& Strict inequality \eqref{eq:report_strict}
& Use ``certified strictly lower Stage-I index'' only on passing years; otherwise report no strict certificate. \\
Rolling M1 supplies an annual upper endpoint
& Complete ledger and concatenation audit
& Full chronology, no overlap/gap, physical replay, nonnegative closing balances
& Use annual scalar bracket only; hourly curves remain diagnostics. \\
M3 is deployable at D--1
& Feature-provenance audit and frozen card
& Every used feature passes 07:00 D--1 and held-year access audit passes
& Use ``D--1 deployable'' only if passed; otherwise describe M3 as retrospective and do not use as primary policy. \\
M3 adds value beyond exposure deformation
& M2, M3, M2-MATCHED
& Positive exposure-effect denominator and positive realized gain after tolerance; shape penalty reported
& State matched-exposure gain only if positive; otherwise state that nominal gain is explained by exposure and/or incurs a shape penalty. \\
All 114 positions support the primary case
& Detailed/reduced validation
& Every 35 multibus position passes \eqref{eq:report_embedding_gate}
& Use 114-position primary only if passed; otherwise 79 direct primary and multibus sensitivity. \\
Retained single-element contingency claim
& Complete contingency registry and independent DC audit
& Every included state has supported one-element semantics, connectivity, ratings, and affine/DC consistency
& Use the exact retained-state phrase; never ``full Dutch N--1.'' \\
Frozen DC evaluator is independently reproduced
& Level-1 status and row table
& Numerical gate, exact IDs, zero false-safe/false-overload beyond tolerance
& State independent DC agreement only if passed; otherwise no row-level policy claim. \\
Selected-state nonlinear detail preserves ordering
& Level-2 AC status and pairwise direction table
& Every non-tied predeclared pair has same DC/AC severity sign
& Use no-reversal wording only if passed; otherwise report reversal frequency and mixed consistency. \\
Background chronology is AC-plausible
& Optional all-hour background AC status
& Declared convergence and feasibility summaries complete; no missing partitions
& Describe background reconstruction only; never annual policy AC feasibility. \\
Common redispatch does not reverse ordering
& Complete redispatch source/input/output/status
& Every non-tied selected pair preserves sign under common action set
& Use only if complete and passed; otherwise remove numerical redispatch claim and retain limitation. \\
Gamma path is stable over the tested range
& Exact gamma and M2-MATCHED case statuses
& Every grid point complete; endpoints reproduce M2/frozen card; no interpolated values
& Describe exact post-freeze path; bootstrap remains descriptive. \\
\bottomrule
\end{longtable}

\section{Limitations and Prohibited Inferences}
\label{sec:report_limitations}

The methodology supports exact statements about the debit geometry and
conditional statements about the declared public-data evaluator. It does not
support the following inferences without separate evidence:
\begin{enumerate}[leftmargin=*,itemsep=3pt]
\item full Dutch N--1 security or operator validation;
\item annual AC-SCOPF or AC-OPF feasibility;
\item welfare, tariff efficiency, customer compensation, actual savings, or
reinforcement deferral;
\item legal equivalence between touched time and physical FDE;
\item uniqueness or price interpretation of the selected Stage-II dual;
\item deployability of realized or scenario residual load at D--1;
\item physical equivalence of the 383-bus and 45-bus layers without the frozen
validation gate;
\item statistical generalization from fixed-policy temporal block resampling;
\item policy stability from global AC/DC correlation alone;
\item redispatch conclusions from incomplete or nonreproducible artifacts.
\end{enumerate}

The hourly testbed is an operational abstraction. The temporal aggregation
result explains which debit properties are preserved and which are not; it
does not recreate unobserved quarter-hour variability. Public network and
scenario data remain approximations. The report therefore requires explicit
source, mapping, baseline-insecurity, reduction-error, and engineering-stress
diagnostics rather than presenting the model as an operator digital twin.

\section{Submission Disposition}
\label{sec:report_disposition}

The scientific specification is internally fixed: one restrictions-only parent
model, one Stage-I/II objective, three debit laws, one allowance-specific M3
fit--validate--freeze sequence, one two-layer Dutch physical case, and one
ordered evidence stack. The report is ready as the methodological and data
companion.

Two factual inputs remain external to the text source supplied for this edition:
(i) the final claim-gated numerical outcome files and (ii) the row-level frozen
bus, element, device, and contingency tables. The document contains exact input
hooks and forbids manual replacement. Submission status becomes final only when
those files are present, hashed, and the claim matrix passes.

\begin{center}
\fbox{\parbox{0.88\textwidth}{
\centering
\textbf{METHOD, DATA SCHEMA, PROOFS, AND EVIDENCE CONTRACT: READY}\\[1mm]
\textbf{FINAL NUMERICAL PROSE AND RAW TOPOLOGY TABLES: PENDING FROZEN ARTIFACTS}
}}
\end{center}

\appendices

\section{Required Metadata Fields}
\label{app:metadata}

\begin{longtable}{@{}p{0.22\textwidth}p{0.70\textwidth}@{}}
\caption{Minimum metadata carried by every scientific object.}\\
\toprule
Object & Required fields \\
\midrule
\endfirsthead
\toprule
Object & Required fields \\
\midrule
\endhead
Time grid & ID, year, timezone, timestamps, interval durations, native resolution, completeness policy, source and hash. \\
Participant & Stable ID, direction, technology/class, province, contracted capacity, detailed and reduced connection IDs, mapping type, balancing-key ID. \\
Allowance & Nominal $\rho$, metric, contract-mask hash, interval count, floor rule, actual vector, units, source, vector hash. \\
Physical case & Network release, bus/element/device IDs, ratings, phase shifts, background state, requests, masks, participant set, contingency set, all matrix hashes. \\
Card & Allowance-specification ID and vector hash, feature schema, provenance hash, model family, preprocessing, parameters, thresholds, rates, bounds, $\lambda_q$, fallback, fit/validation years, selection values, code and source hashes. \\
Experiment & Case ID, physical-case ID, time-grid ID, rule, allowance source/hash, card and feature IDs, look-ahead, implementation block, information level, solver profile, code hash. \\
Solution & Stage certificates, row closure, primal actions, duals where applicable, objective values, participant and hourly metrics, replay status, solver environment, artifact hashes. \\
Engineering case & Selection reason, year, timestamp, policy, state, injection hash, row and branch maps, P/Q assumption, solver attempt sequence, result status. \\
\bottomrule
\end{longtable}

\section{Required Unit and Regression Tests}
\label{app:tests}

Before any annual campaign, the implementation must pass at least the following
focused tests:
\begin{enumerate}[leftmargin=*,itemsep=2pt]
\item explicit duration changes every objective, debit, ledger, and metric by the
expected factor;
\item zero request leaves the contract-duration allowance unchanged;
\item floor rounding differs correctly from ceiling in fractional cases;
\item M1 links reject positive submaterial actions and accept zero/full actions;
\item M1-LP projection agrees with an explicit relaxed-$u$ formulation on small
instances;
\item the source-of-separation identity and set hierarchy hold on random small
instances;
\item Stage-II duals are unchanged by the absence of Stage III and are never
extracted after the selector;
\item sparse M3 directions freeze exactly at $q=1$;
\item M2-MATCHED contains the completed M3 trajectory and uses the exact vector
hash;
\item temporal aggregation reproduces M2 and the exposure-weighted M3 identity
and exhibits the M1 counterexample;
\item safe screening never removes a row that can violate on enumerated small
boxes;
\item row closure and independent replay detect deliberately hidden violations;
\item background and participant transactions balance at every interval;
\item $\bm G^c=\bm H^c\bm A$ holds with exact row and participant IDs;
\item phase-shift test cases require and reproduce nonzero $\bm h^c$;
\item parallel and islanding contingency cases receive the correct labels;
\item held-out files are inaccessible to fit/selection code before the card
freeze;
\item AC row maps detect reversed orientation and MW/MVA misuse;
\item all completed status files fail when a tracked source or artifact hash is
changed.
\end{enumerate}

\section{Implementation Dependency Order}
\label{app:dependency}

The required implementation order is:
\begin{enumerate}[leftmargin=*,itemsep=1pt]
\item methodology freeze and machine-readable contract;
\item experiment, time-grid, allowance, physical-case, and solution schemas;
\item input loader and exact identifiers;
\item common Stage-I/II/III solver, screening, closure, and replay;
\item theorem-level and solver-level unit tests;
\item M3 teacher, provenance gate, features, predictors, card QP, and freeze;
\item Dutch source extraction and component schemas;
\item technology allocations and dispatch/boundary state builder;
\item contingency and affine evaluator builder;
\item detailed/reduced participant effects and embedding validation;
\item 2018 reference extraction and plausibility audit;
\item immutable physical-case assembly;
\item annual, rolling, matched-exposure, control, and gamma runners;
\item independent DC, selected-state AC, background AC, and redispatch audits;
\item claim-to-artifact evaluator and reporting scripts;
\item final numerical campaign and only then the numerical clauses in the abstract,
results, and conclusion.
\end{enumerate}
This order prevents a card from being frozen against a changing objective or
physical case and prevents manuscript claims from preceding evidence.

\begin{thebibliography}{9}
\bibitem{Zomerdijk2022}
W.~Zomerdijk, D.~Gusain, P.~Palensky, and M.~Cvetkovic,
``Open data based model of the Dutch high-voltage power system,'' in
\emph{Proc. IEEE PES ISGT Europe}, 2022, pp.~1--6.
\bibitem{ZomerdijkData}
W.~Zomerdijk, ``Open Data Based Model of the Dutch High-Voltage Power System,''
dataset and software, Zenodo, 2025, doi: 10.5281/zenodo.15223447.
\bibitem{Stott2009}
B.~Stott, J.~Jardim, and O.~Alsac, ``DC power flow revisited,''
\emph{IEEE Trans. Power Syst.}, vol.~24, no.~3, pp.~1290--1300, 2009.
\bibitem{Thurner2018}
L.~Thurner \emph{et al.}, ``pandapower---An open-source Python tool for
convenient modeling, analysis, and optimization of electric power systems,''
\emph{IEEE Trans. Power Syst.}, vol.~33, no.~6, pp.~6510--6521, 2018.
\bibitem{BonnansShapiro2000}
J.~F. Bonnans and A.~Shapiro, \emph{Perturbation Analysis of Optimization
Problems}. New York, NY, USA: Springer, 2000.
\bibitem{Kunsch1989}
H.~R. K{\"u}nsch, ``The jackknife and the bootstrap for general stationary
observations,'' \emph{Ann. Statist.}, vol.~17, no.~3, pp.~1217--1241, 1989.
\bibitem{PolitisRomano1994}
D.~N. Politis and J.~P. Romano, ``The stationary bootstrap,''
\emph{J. Amer. Statist. Assoc.}, vol.~89, no.~428, pp.~1303--1313, 1994.
\end{thebibliography}
\end{document}
